\section{Results}

\subsection{Sample construction}
The source contained \InitialWindows{} team--match five-minute windows from \MatchCount{} matches and \TeamCount{} teams. After the common score-state, period, and numerical-state restrictions and the opponent-pass threshold, \IntensityWindows{} windows entered the intensity model. The sequence denominator restrictions retained \EfficiencyWindows{} windows for the efficiency model. The stepwise exclusion table reconciles every reduction; no match or team was silently removed. Drawing contributed the largest exposure total, while leading and trailing exposure seconds were symmetric in the fixed-window construction. Opportunity totals by state are shown in Table~1.

\subsection{Descriptive results}
Pooled Pressure events per 30 opponent passes were \DescriptiveIntensityLeadingPooled{} when leading, \DescriptiveIntensityDrawingPooled{} when drawing, and \DescriptiveIntensityTrailingPooled{} when trailing. The corresponding unweighted mean window rates were \DescriptiveIntensityLeadingMean{}, \DescriptiveIntensityDrawingMean{}, and \DescriptiveIntensityTrailingMean{}. These differ because pooled rates weight windows by opponent-pass exposure, whereas unweighted summaries give each eligible window equal weight.

Pooled five-second sequence-regain proportions were \DescriptiveEfficiencyLeadingPooled{}, \DescriptiveEfficiencyDrawingPooled{}, and \DescriptiveEfficiencyTrailingPooled{} for leading, drawing, and trailing. Unweighted mean window proportions were \DescriptiveEfficiencyLeadingMean{}, \DescriptiveEfficiencyDrawingMean{}, and \DescriptiveEfficiencyTrailingMean{}. These are unadjusted summaries and should not be read as model-standardized comparisons.

\subsection{Primary pressing-intensity model}
Relative to drawing, trailing was associated with a higher adjusted Pressure-event rate: IRR \IntensityTrailingEffect{} (log-rate coefficient \IntensityTrailingCoef{}, 95\% CI \IntensityTrailingCILow{} to \IntensityTrailingCIHigh{}). Predicted Pressure events per 30 opponent passes were \IntensityTrailingPrediction{} when trailing and \IntensityDrawingPrediction{} when drawing, an adjusted difference of approximately \IntensityTrailingDifference\%. Leading versus drawing had IRR \IntensityLeadingEffect{} (coefficient \IntensityLeadingCoef{}, 95\% CI \IntensityLeadingCILow{} to \IntensityLeadingCIHigh{}); predicted rates were \IntensityLeadingPrediction{} and \IntensityDrawingPrediction{}, respectively, an adjusted difference of approximately \IntensityLeadingDifference\%.

\begin{figure}[ht]
\centering
\includegraphics[width=.82\linewidth]{../../outputs/figures/pressing_score_state/figure_3.pdf}
\caption{Adjusted pressing intensity. Predicted provider Pressure events per 30 opponent passes from the confirmatory Poisson model, using team--match five-minute windows and drawing as reference. Team fixed effects and match-clustered uncertainty are used; mixed state, numerical inequality, extra time, and windows with fewer than five opponent passes are excluded. Exact confidence limits are preserved in the source CSV.}
\label{fig:intensity}
\end{figure}

\subsection{Primary regain-efficiency model}
Trailing versus drawing had an estimated sequence-regain odds ratio of \EfficiencyTrailingEffect{} (log-odds coefficient \EfficiencyTrailingCoef{}, 95\% CI \EfficiencyTrailingCILow{} to \EfficiencyTrailingCIHigh{}). Adjusted predicted probabilities were \EfficiencyTrailingPrediction{} and \EfficiencyDrawingPrediction{}, a \EfficiencyTrailingDifference{} percentage-point difference. Leading versus drawing had OR \EfficiencyLeadingEffect{} (coefficient \EfficiencyLeadingCoef{}, 95\% CI \EfficiencyLeadingCILow{} to \EfficiencyLeadingCIHigh{}); predicted probabilities were \EfficiencyLeadingPrediction{} and \EfficiencyDrawingPrediction{}, a \EfficiencyLeadingDifference{} percentage-point difference.

\begin{figure}[ht]
\centering
\includegraphics[width=.82\linewidth]{../../outputs/figures/pressing_score_state/figure_4.pdf}
\caption{Adjusted five-second Pressure-sequence regain efficiency. Predicted probabilities from the confirmatory grouped-binomial model, using team--match five-minute windows and drawing as reference. Team fixed effects and match-clustered uncertainty are used; mixed state, numerical inequality, extra time, and windows with fewer than three sequences are excluded.}
\label{fig:efficiency}
\end{figure}

\subsection{Multiplicity-adjusted inference}
Holm-adjusted primary-family values were \IntensityTrailingHolmP{} for the trailing intensity contrast and \EfficiencyTrailingHolmP{} for trailing efficiency. The leading values were \IntensityLeadingHolmP{} and \EfficiencyLeadingHolmP{}. Interpretation rests on transformed estimates, adjusted predictions, and interval width in addition to these values.

\begin{table}[ht]
\centering
\small
\caption{Primary adjusted models. Team--match five-minute windows; drawing reference; team fixed effects and match-clustered standard errors; mixed states, numerical inequality, extra time, and unreliable denominators excluded. Confirmatory analyses.}
\begin{tabular}{lllll}
\hline
Outcome & Contrast & Estimate & 95\% CI & Transformed \\
pressure\_events & score\_state[leading] & -0.032 & [-0.103, 0.038] & 0.968 \\
pressure\_events & score\_state[trailing] & 0.135 & [0.042, 0.229] & 1.145 \\
sequence\_regains\_5s/pressure\_sequences & score\_state[leading] & -0.086 & [-0.231, 0.060] & 0.918 \\
sequence\_regains\_5s/pressure\_sequences & score\_state[trailing] & 0.231 & [0.080, 0.382] & 1.260 \\
\hline
\end{tabular}
\end{table}


\subsection{Secondary and exploratory results}
High-pressure share and event-level regain models supplied secondary checks on location and measurement unit. Their estimated trailing contrasts were positive, while the event and sequence estimands remained distinct. Provider counterpress share, high-regain rate, Pressure height, sequence frequency, and classic defensive-action counts are retained in Table~3 rather than selected by statistical appearance. Post-regain shot production is secondary; xG and xT per successful regain are exploratory. Same-possession goal is explicitly unavailable, and provider-augmented PPDA remains an unavailable sensitivity entry in the inferential table.

\begin{table}[ht]
\centering
\small
\caption{Secondary and exploratory adjusted analyses. Team--match five-minute windows; drawing reference; match-clustered uncertainty and outcome-specific exclusions. These analyses are not confirmatory primary outcomes.}
\begin{tabular}{lllll}
\hline
Model & Contrast & Estimate & 95\% CI & Role \\
secondary\_high\_pressure\_share & score\_state[leading] & -0.170 & [-0.326, -0.014] & secondary \\
secondary\_high\_pressure\_share & score\_state[trailing] & 0.266 & [0.117, 0.415] & secondary \\
secondary\_event\_regain\_5s & score\_state[leading] & -0.123 & [-0.278, 0.032] & secondary \\
secondary\_event\_regain\_5s & score\_state[trailing] & 0.215 & [0.064, 0.367] & secondary \\
secondary\_counterpress\_share & score\_state[leading] & -0.163 & [-0.317, -0.010] & secondary \\
secondary\_counterpress\_share & score\_state[trailing] & 0.103 & [-0.054, 0.260] & secondary \\
secondary\_high\_regain\_rate & score\_state[leading] & 0.012 & [-0.226, 0.249] & secondary \\
secondary\_high\_regain\_rate & score\_state[trailing] & 0.104 & [-0.162, 0.370] & secondary \\
secondary\_post\_regain\_shot\_10s & score\_state[leading] & 0.170 & [-0.288, 0.628] & exploratory \\
secondary\_post\_regain\_shot\_10s & score\_state[trailing] & -0.212 & [-0.737, 0.312] & exploratory \\
secondary\_pressure\_sequence\_frequency & score\_state[leading] & 0.080 & [-0.005, 0.164] & secondary \\
secondary\_pressure\_sequence\_frequency & score\_state[trailing] & -0.011 & [-0.097, 0.076] & secondary \\
secondary\_classic\_defensive\_action\_rate & score\_state[leading] & -0.011 & [-0.189, 0.167] & secondary \\
secondary\_classic\_defensive\_action\_rate & score\_state[trailing] & -0.005 & [-0.173, 0.163] & secondary \\
exploratory\_post\_regain\_xg\_per\_success & score\_state[leading] & 0.001 & [-0.004, 0.006] & exploratory \\
exploratory\_post\_regain\_xg\_per\_success & score\_state[trailing] & -0.001 & [-0.007, 0.005] & exploratory \\
exploratory\_post\_regain\_xt\_per\_success & score\_state[leading] & -0.006 & [-0.012, 0.001] & exploratory \\
exploratory\_post\_regain\_xt\_per\_success & score\_state[trailing] & 0.004 & [-0.003, 0.010] & exploratory \\
secondary\_mean\_pressure\_height & score\_state[leading] & -1.285 & [-3.387, 0.817] & secondary \\
secondary\_mean\_pressure\_height & score\_state[trailing] & 4.148 & [2.264, 6.031] & secondary \\
\hline
\end{tabular}
\end{table}


\begin{figure}[ht]
\centering
\includegraphics[width=.82\linewidth]{../../outputs/figures/pressing_score_state/figure_5.pdf}
\caption{Intensity--efficiency comparison. Trailing-minus-drawing adjusted contrasts are shown on interpretable but different scales: percentage change in Pressure events per opponent-pass exposure and percentage-point change in sequence-regain probability. Team--match five-minute windows, drawing reference, primary exclusions, team fixed effects, and match-clustered uncertainty apply. Both outcomes are confirmatory, but their scales are not interchangeable.}
\end{figure}
